%%%%%%%%%%%%%%%%%%%%%%%%%%%%%%%%%%%%%%%%%
% Classicthesis Typographic Thesis
% LaTeX Template
% Version 1.4 (1/1/16)
%
% This template has been downloaded from:
% http://www.LaTeXTemplates.com
%
% Original author:
% 
% License:
% GNU General Public License (v2)
%
% General Tips:
% 1) Make sure to edit the classicthesis-config.file
% 2) New enumeration (A., B., C., etc in small caps): \begin{aenumerate} \end{aenumerate}
% 3) For margin notes: \marginpar or \graffito{}
% 4) Do not use bold fonts in this style, it is designed around them
% 5) Use tables as in the examples
% 6) See classicthesis-preamble.sty for useful commands
%
%%%%%%%%%%%%%%%%%%%%%%%%%%%%%%%%%%%%%%%%%

%----------------------------------------------------------------------------------------
%	PACKAGES AND OTHER DOCUMENT CONFIGURATIONS
%----------------------------------------------------------------------------------------

\documentclass[
		%headinclude=true,
		%twoside,
		openright,
		titlepage,
		numbers=noenddot,
		headinclude,
		%1headlines,
	 	%footinclude=true,
	 	cleardoublepage=empty,
		dottedtoc, % Make page numbers in the table of contents flushed right with dots leading to them
		%BCOR=5mm,
		paper=a4,
		fontsize=11pt, % Binding correction, paper type and font size
		french,
		american, % Languages, change this to your language(s)
		]{scrreprt} 
                
% Includes the file which contains all the document configurations and packages - make sure to edit this file
%%%%%%%%%%%%%%%%%%%%%%%%%%%%%%%%%%%%%%%%%
% Classicthesis Typographic Thesis
% Configuration File
%
% This file has been downloaded from:
% http://www.LaTeXTemplates.com
%
% Original author:
% André Miede (http://www.miede.de) with extensive commenting changes by:
% Vel (vel@LaTeXTemplates.com)
%
% License:
% GNU General Public License (v2)
%
% Important note:
% The main lines to change in this file are in the DOCUMENT VARIABLES
% section, the rest of the file is for advanced configuration.
%
%%%%%%%%%%%%%%%%%%%%%%%%%%%%%%%%%%%%%%%%%

%----------------------------------------------------------------------------------------
%	CHARACTER ENCODING
%----------------------------------------------------------------------------------------

\PassOptionsToPackage{utf8}{inputenc} % Set the encoding of your files. UTF-8 is the only sensible encoding nowadays. If you can't read äöüßáéçèê∂åëæƒÏ€ then change the encoding setting in your editor, not the line below. If your editor does not support utf8 use another editor!
\usepackage{inputenc}

%----------------------------------------------------------------------------------------
%	DOCUMENT VARIABLES
%	Fill in the lines below to enter your information into the thesis template
%	Each of the commands can be cited anywhere in the thesis
%----------------------------------------------------------------------------------------

% Remove drafting to get rid of the '[ Date - classicthesis version 4.0 ]' text at the bottom of every page
\PassOptionsToPackage{eulerchapternumbers,listings,drafting, pdfspacing, subfig,beramono,eulermath,parts}{classicthesis}
% Available options: drafting parts nochapters linedheaders eulerchapternumbers beramono eulermath pdfspacing minionprospacing tocaligned dottedtoc manychapters listings floatperchapter subfig

%\newcommand{\myTitle}{Apprentissage supervisé, la reconnaissance des plaques d'immatriculation utilisant l'algorithme de descente de gradient stochastique\xspace}
\newcommand{\myTitle}{Machine Learning: Étude de la minimisation d'erreur dans l'apprentissage supervisé, avec une application de la technologie ANPR\xspace}
\newcommand{\myTitleEn}{Supervised learning, License Recognition using the stochastic gradient descent algorithm\xspace}
\newcommand{\mySubtitle}{TFC\xspace}
\newcommand{\myDegree}{Licencié en informatique\xspace}
\newcommand{\myName}{\textsc{TSHELEKA KAJILA} Hassan\xspace}
%\newcommand{\myProf}{Prof. BAGULA Antoine\xspace}
\newcommand{\myProf}{Prof. MASAKUNA Jordan\xspace}
\newcommand{\myOtherProf}{Put name here\xspace}
\newcommand{\mySupervisor}{Mr. \textsc{MWAHU} Auguste\xspace}
\newcommand{\myFaculty}{Sciences Informatiques\xspace}
\newcommand{\myDepartment}{\href{https://www.unhorizons.org}{Calcul Scientifique}\xspace}
\newcommand{\myUni}{\href{http://www.unhorizons.org}{Université Nouveaux Horizons}\xspace}
\newcommand{\myLocation}{Lubumbashi\xspace}
\newcommand{\myTime}{Janvier 2022\xspace}
\newcommand{\myVersion}{version 4.2\xspace}

%----------------------------------------------------------------------------------------
%	USEFUL COMMANDS
%----------------------------------------------------------------------------------------

\newcommand{\ie}{i.\,e.}
\newcommand{\Ie}{I.\,e.}
\newcommand{\eg}{e.\,g.}
\newcommand{\Eg}{E.\,g.} 
\newcommand{\cf}{\emph{cf.}} 

\newcounter{dummy} % Necessary for correct hyperlinks (to index, bib, etc.)
\providecommand{\mLyX}{L\kern-.1667em\lower.25em\hbox{Y}\kern-.125emX\@}
\newlength{\abcd} % for ab..z string length calculation

%----------------------------------------------------------------------------------------
%	PACKAGES
%----------------------------------------------------------------------------------------

\usepackage{lipsum} % Used for inserting dummy 'Lorem ipsum' text into the template

%------------------------------------------------

%\PassOptionsToPackage{ngerman,american}{babel}  % Change this to your language(s)
% Spanish languages need extra options in order to work with this template
%\PassOptionsToPackage{spanish,es-lcroman}{babel}
\usepackage{babel}

%------------------------------------------------			

\usepackage{csquotes}
\PassOptionsToPackage{%
%backend=biber, % Instead of bibtex
backend=bibtex8,bibencoding=ascii,%
language=auto,%
style=numeric-comp,%
%style=authoryear-comp, % Author 1999, 2010
%bibstyle=authoryear,dashed=false, % dashed: substitute rep. author with ---
sorting=nyt, % name, year, title
maxbibnames=10, % default: 3, et al.
%backref=true,%
natbib=true % natbib compatibility mode (\citep and \citet still work)
}{biblatex}
\usepackage{biblatex}
 
 %------------------------------------------------

\PassOptionsToPackage{fleqn}{amsmath} % Math environments and more by the AMS 
 \usepackage{amsmath}
 
 %------------------------------------------------

\PassOptionsToPackage{T1}{fontenc} % T2A for cyrillics
\usepackage{fontenc}

%------------------------------------------------

\usepackage{textcomp} % Fix warning with missing font shapes

%------------------------------------------------

\usepackage{scrhack} % Fix warnings when using KOMA with listings package  

%------------------------------------------------

\usepackage{xspace} % To get the spacing after macros right

%------------------------------------------------

\usepackage{mparhack} % To get marginpar right

%------------------------------------------------

\usepackage{fixltx2e} % Fixes some LaTeX stuff 

%------------------------------------------------

\PassOptionsToPackage{smaller}{acronym} % Include printonlyused in the first bracket to only show acronyms used in the text
\usepackage{acronym} % Nice macros for handling all acronyms in the thesis

%\renewcommand*{\acsfont}[1]{\textssc{#1}} % For MinionPro
\renewcommand*{\aclabelfont}[1]{\acsfont{#1}}

%------------------------------------------------

\PassOptionsToPackage{pdftex}{graphicx}
\usepackage{graphicx} 

%----------------------------------------------------------------------------------------
%	FLOATS: TABLES, FIGURES AND CAPTIONS SETUP
%----------------------------------------------------------------------------------------

\usepackage{tabularx} % Better tables
\setlength{\extrarowheight}{3pt} % Increase table row height
\newcommand{\tableheadline}[1]{\multicolumn{1}{c}{\spacedlowsmallcaps{#1}}}
\newcommand{\myfloatalign}{\centering} % To be used with each float for alignment
\usepackage{caption}
\captionsetup{font=small}
\usepackage{subfig}  

%----------------------------------------------------------------------------------------
%	CODE LISTINGS SETUP
%----------------------------------------------------------------------------------------

\usepackage{listings} 
%\lstset{emph={trueIndex,root},emphstyle=\color{BlueViolet}}%\underbar} % For special keywords
\lstset{language=[LaTeX]Tex,%C++ % Specify the language(s) for listings here
	morekeywords={PassOptionsToPackage,selectlanguage},
	keywordstyle=\color{RoyalBlue}, % Add \bfseries for bold
	basicstyle=\small\ttfamily, % Makes listings a smaller font size and a different font
	%identifierstyle=\color{NavyBlue}, % Color of text inside brackets
	commentstyle=\color{Green}\ttfamily, % Color of comments
	stringstyle=\rmfamily, % Font type to use for strings
	numbers=left, % Change left to none to remove line numbers
	numberstyle=\scriptsize, % Font size of the line numbers
	stepnumber=5, % Increment of line numbers
	numbersep=8pt, % Distance of line numbers from code listing
	showstringspaces=false, % Sets whether spaces in strings should appear underlined
	breaklines=true, % Force the code to stay in the confines of the listing box
	%frameround=ftff, % Uncomment for rounded frame
	%frame=single, % Frame border - none/leftline/topline/bottomline/lines/single/shadowbox/L
	belowcaptionskip=.75\baselineskip % Space after the "Listing #: Desciption" text and the listing box
}

%----------------------------------------------------------------------------------------
%	HYPERREFERENCES
%----------------------------------------------------------------------------------------

\PassOptionsToPackage{pdftex,hyperfootnotes=false,pdfpagelabels}{hyperref}
\usepackage{hyperref}  % backref linktocpage pagebackref
\pdfcompresslevel=9
\pdfadjustspacing=1

\hypersetup{
	% Uncomment the line below to remove all links (to references, figures, tables, etc), useful for b/w printouts
	%draft, 
	colorlinks=true, linktocpage=true, pdfstartpage=3, pdfstartview=FitV,
	% Uncomment the line below if you want to have black links (e.g. for printing black and white)
	%colorlinks=false, linktocpage=false, pdfborder={0 0 0}, pdfstartpage=3, pdfstartview=FitV, 
	breaklinks=true, pdfpagemode=UseNone, pageanchor=true, pdfpagemode=UseOutlines,%
	plainpages=false, bookmarksnumbered, bookmarksopen=true, bookmarksopenlevel=1,%
	hypertexnames=true, pdfhighlight=/O,%nesting=true,%frenchlinks,%
	urlcolor=webbrown, linkcolor=RoyalBlue, citecolor=webgreen, %pagecolor=RoyalBlue,%
	    %urlcolor=Black, linkcolor=Black, citecolor=Black, %pagecolor=Black,%
	%------------------------------------------------
	% PDF file meta-information
	pdftitle={\myTitle},
	pdfauthor={\textcopyright\ \myName, \myUni, \myFaculty},
	pdfsubject={thesis},
	pdfkeywords={ECW},
	pdfcreator={},
	pdfproducer={LaTeX}
	%------------------------------------------------
}

%----------------------------------------------------------------------------------------
%	AUTOREFERENCES SETUP
%	Redefines how references in text are prefaced for different 
%	languages (e.g. "Section 1.2" or "section 1.2")
%----------------------------------------------------------------------------------------

\makeatletter
\@ifpackageloaded{babel}
{
\addto\extrasamerican{
\renewcommand*{\figureautorefname}{Figure}
\renewcommand*{\tableautorefname}{Table}
\renewcommand*{\partautorefname}{Part}
\renewcommand*{\chapterautorefname}{Chapter}
\renewcommand*{\sectionautorefname}{Section}
\renewcommand*{\subsectionautorefname}{Section}
\renewcommand*{\subsubsectionautorefname}{Section}
}
\addto\extrasngerman{
\renewcommand*{\paragraphautorefname}{Absatz}
\renewcommand*{\subparagraphautorefname}{Unterabsatz}
\renewcommand*{\footnoteautorefname}{Fu\"snote}
\renewcommand*{\FancyVerbLineautorefname}{Zeile}
\renewcommand*{\theoremautorefname}{Theorem}
\renewcommand*{\appendixautorefname}{Anhang}
\renewcommand*{\equationautorefname}{Gleichung}
\renewcommand*{\itemautorefname}{Punkt}
}
\providecommand{\subfigureautorefname}{\figureautorefname} % Fix to getting autorefs for subfigures right
}{\relax}
\makeatother

%----------------------------------------------------------------------------------------

\usepackage{classicthesis} 

%----------------------------------------------------------------------------------------
%	CHANGING TEXT AREA 
%----------------------------------------------------------------------------------------

%\linespread{1.05} % a bit more for Palatino
%\areaset[current]{312pt}{761pt} % 686 (factor 2.2) + 33 head + 42 head \the\footskip
%\setlength{\marginparwidth}{7em}%
%\setlength{\marginparsep}{2em}%

%----------------------------------------------------------------------------------------
%	USING DIFFERENT FONTS
%----------------------------------------------------------------------------------------

%\usepackage[oldstylenums]{kpfonts} % oldstyle notextcomp
%\usepackage[osf]{libertine}
%\usepackage[light,condensed,math]{iwona}
%\renewcommand{\sfdefault}{iwona}
%\usepackage{lmodern} % <-- no osf support :-(
%\usepackage{cfr-lm} % 
%\usepackage[urw-garamond]{mathdesign} <-- no osf support :-(
%\usepackage[default,osfigures]{opensans} % scale=0.95 
%\usepackage[sfdefault]{FiraSans}
%\usepackage{lipsum}
\usepackage{float}
\usepackage{pdfpages}
\usepackage{wrapfig}


\addbibresource{mybib.bib} % The file housing your bibliography
%\addbibresource[label=ownpubs]{Self_Publications.bib} % Uncomment for optional self-publications

%\hyphenation{Put special hyphenation here}

\areaset{15cm}{25.8cm}

\parindent 0pt
\parskip 6pt

\begin{document}

%\renewcommand{\thesubsubsection}{\alph{subsection}}
\renewcommand{\thesubsubsection}{\alph{subsubsection}.}
\setcounter{chapter}{-1}


\frenchspacing % Reduces space after periods to make text more compact

\raggedbottom % Makes all pages the height of the text on that page

\selectlanguage{french} % Select your default language - e.g. american or ngerman

%\renewcommand*{\bibname}{new name} % Uncomment to change the name of the bibliography
%\setbibpreamble{} % Uncomment to include a preamble to the bibliography - some text before the reference list starts

\pagenumbering{roman} % Roman page numbering prior to the start of the thesis content (i, ii, iii, etc)

\pagestyle{plain} % Suppress headers for the pre-content pages

%----------------------------------------------------------------------------------------
%	PRE-CONTENT THESIS PAGES
%----------------------------------------------------------------------------------------

% Title Page

\begin{titlepage}

%\begin{addmargin}[-1cm]{-3cm}
\begin{center}
\large

\vspace*{.03\textheight}


 

%\text{word or phrase}
%\vspace{1.5cm} 
\includegraphics[width=12cm]{images/logo_unh} \\[0.75cm]

{\LARGE{{\textsc{FACULT\'E DES SCIENCES INFORMATIQUES}}} }\\[0.2cm]

{\LARGE{{{Calcul Scientifique}}} }\\
\medskip % Picture

\vspace{2.5cm} 

%{\hrule width5cm}  % Horizontal line
\begin{center}
	\rule{0.9\textwidth}{.2pt}
\end{center}
\vspace{0.2cm}
{\LARGE \bfseries {\myTitle}}%\vspace{0.4cm} % Thesis title
%{\hrule width5cm} 
\begin{center}
	\rule{0.9\textwidth}{.2pt}
\end{center}

%\color{Maroon}\spacedallcaps{\myTitle} \\ \bigskip % Thesis title

\textit{Travail de fin d’études présenté  et défendu en vue de l’obtention du grade de licencié en \myDepartment.}\\[2cm]

\begin{minipage}[t]{0.4\textwidth}
	\begin{flushleft} \large
		%\emph{Auteur:}\\
		%{\myName} % Author name - remove the \href bracket to remove the link
	\end{flushleft}
\end{minipage}
\begin{minipage}[t]{0.5\textwidth}
	\begin{flushright} \large
		\emph{Auteur:}
		\textbf{{\myName}}\\
		\emph{Directeur:} 
		%\href{https://cs.uwc.ac.za/anttoine-bagula/}{\myProf} % Supervisor name - remove the \href bracket to remove the link
		\textbf{{\myProf}}  
	\end{flushright}
\end{minipage}\\[3cm]

%\spacedlowsmallcaps{\myName} % Your name

\vfill


% Travail de fin d’études présenté en vue de l’obtention du grade de licence en Sciences Informatiques
%Mémoire présenté à la Faculté des \myFaculty en vue de l'obtention du grade de \myDegree.
%Travail de fin d’études présenté en vue de l’obtention du grade de \myDegree en Sciences Informatiques.

 % University requirement text
%\textit{département}\\[0.2cm]
%\myDepartment \\[0.75cm]
%\medskip % Thesis subtitle
%\myDegree \\
%\myDepartment \\
%\myFaculty \\
%\myUni \\ \bigskip

{\large \today}\\

%\myTime\ - % Time and version



\end{center}
%\end{addmargin}

\end{titlepage} % Main title page
%%\includegraphics[]{gfx/garde} 

%\includepdf[pages={1}]{gfx/garde.pdf}

% Back of the title page

\thispagestyle{empty}

\hfill

\vfill

\noindent\myName: \textit{\myTitle,} \mySubtitle, %\myDegree, 
\textcopyright\ \myTime

% You may wish to do something with the back of the title page, such as including your supervisors, location or time frame of the work. Below is an example of doing so although you may want to tweak it to your liking.

%\bigskip

%\noindent\spacedlowsmallcaps{Supervisors}: \\
%\myProf \\
%\myOtherProf \\ 
%\mySupervisor

%\medskip \\

%\noindent\spacedlowsmallcaps{Location}: \\
%\myLocation

%\medskip \\

%\noindent\spacedlowsmallcaps{Time Frame}: \\
%\myTime
 % Back of the title page

%\cleardoublepage% Dedication

\thispagestyle{empty}
\refstepcounter{dummy}

\pdfbookmark[1]{Dedication}{Dedication} % Bookmark name visible in a PDF viewer

\vspace*{3cm}

\begin{center}
\emph{Ohana} means family. \\
Family means nobody gets left behind, or forgotten. \\ \medskip
--- Lilo \& Stitch    
\end{center}

\medskip

\begin{center}
Dedicated to the loving memory of Rudolf Miede. \\ \smallskip
1939\,--\,2005
\end{center} % Dedication page

%\cleardoublepage\include{FrontBackMatter/Foreword} % Uncomment and create a Foreword.tex to include a foreword

%\cleardoublepage% Abstract

\renewcommand{\abstractname}{Résumé} % Uncomment to change the name of the abstract

\pdfbookmark[1]{Résumé}{Résumé} % Bookmark name visible in a PDF viewer
\addcontentsline{toc}{chapter}{Résumé}
\begingroup
\let\clearpage\relax
\let\cleardoublepage\relax
\let\cleardoublepage\relax

\chapter*{Résumé}
	Au cours de la dernière décennie, la taille des données a augmenté plus rapidement que la vitesse des processeurs. Dans ce contexte, faire un traitement de {reconnaissance} des formes dans des images et vidéos, les ensembles de données d'entraînement pour les problèmes de détection d'objets sont généralement très volumineux et les capacités des méthodes d'apprentissage automatique statistique sont limitées par le temps de calcul plutôt que par la taille de l'échantillon. 
	
	Le cas des problèmes d'apprentissage à grande échelle implique la complexité de calcul de l'algorithme d'optimisation sous-jacent de manière non triviale. Des algorithmes d'optimisation improbables tels que la \textbf{descente de gradient stochastique} (en anglais: \textbf{Stochastic Gradient Descent} ou SGD) montre des performances étonnantes pour les problèmes à grande échelle, lorsque l'ensemble d'apprentissage est volumineux. \\
	En particulier, les variants du SGD n'utilisent qu'un seul nouvel échantillon d'apprentissage à chaque itération, sont asymptotiquement efficaces après un seul passage sur l'ensemble d'apprentissage.	
	
	Ce travail vise à proposer une méthode  intelligente, basée sur l'intelligence artificielle, qui permet aux ordinateurs et aux systèmes informatiques de dériver des informations significatives à partir d'images numériques, de vidéos et d'autres entrées visuelles, avec un coût plus bas que possible. Dans notre contexte la reconnaissance des plaques d’immatriculation des véhicules à l'aide d’un classificateur de la famille de descente de gradient stochastique. Pour minimiser la \textbf{fonction coût} du classificateur, la SGD adopte un modèle d'optimisation convexe. De plus, pour augmenter la vitesse de convergence du classificateur, la descente de gradient stochastique, à chaque étape, elle tire un échantillon aléatoire de l'ensemble des fonctions ($f_i$), de la fonction objectif, constituant la somme.
	 
	\begin{center}
		
	%\url{https://plg.uwaterloo.ca/~migod/research/beckOOPSLA.html}
	\end{center}
	\textbf{Mots clés~:} Apprentissage supervisé, vision par ordinateur, Descente de gradient stochastique, Adaline, ALPR. 

\endgroup			

\vfill

\pagebreak



\pdfbookmark[2]{Abstract}{Abstract} % Bookmark name visible in a PDF viewer

\begingroup
\let\clearpage\relax
\let\cleardoublepage\relax
\let\cleardoublepage\relax

\chapter*{Abstract}
	Over the past decade, data size has grown faster than processor speeds. In this context, doing pattern recognition processing in real-time videos, training datasets for object detection problems are usually very large, and the capabilities of statistical machine learning methods are limited by computation time rather than sample size.
	
	The case of large scale learning problems involves the computational complexity of the underlying optimization algorithm in a nontrivial way.\\
	Improbable optimization algorithms such as \textbf{Stochastic Gradient Descent} (SGD) show amazing performance for large scale problems, when the training set is bulky.\\
	In particular, SGD variants use only one new training sample at each iteration, are asymptotically efficient after a single pass over the training set.
	
	This work aims to provide an intelligent method, based on artificial intelligence, that allows computers and computer systems to derive meaningful information from digital images, videos and other visual inputs, with a lower cost. as possible. In our context the recognition of vehicle license plates using a classifier of the family of stochastic gradient descent. To minimize the \textbf{cost function} of the classifier, the SGD adopts a convex optimization model. Moreover, to increase the speed of convergence of the classifier, the stochastic gradient descent, at each step, it draws a random sample from the set of functions ($f_i$), of the objective function, constituting the sum.
	
	\begin{center}
		
		%\url{https://plg.uwaterloo.ca/~migod/research/beckOOPSLA.html}
	\end{center}
	\textbf{Key words~:} Supervised learning, computer vision, Stochastic gradient descent, Adaline, ALPR.

	

\endgroup			

\vfill % Abstract page

%\cleardoublepage% Publications - a page listing research articles written using content in the thesis

\pdfbookmark[1]{Publications}{Publications} % Bookmark name visible in a PDF viewer

\chapter*{Publications} % Publications page text

Some ideas and figures have appeared previously in the following publications:\\

\noindent Put your publications from the thesis here. The packages \texttt{multibib} or \texttt{bibtopic} etc. can be used to handle multiple different bibliographies in your document.

%\begin{refsection}[ownpubs]
%    \small
%    \nocite{*} % is local to to the enclosing refsection
%    \printbibliography[heading=none]
%\end{refsection}

%\emph{Attention}: This requires a separate run of \texttt{bibtex} for your \texttt{refsection}, \eg, \texttt{ClassicThesis1-blx} for this file. You might also use \texttt{biber} as the backend for \texttt{biblatex}. See also \url{http://tex.stackexchange.com/questions/128196/problem-with-refsection}. % Publications from the thesis page

\cleardoublepage% Acknowledgements

\pdfbookmark[1]{Acknowledgements}{Acknowledgements} % Bookmark name visible in a PDF viewer

\begin{flushright}{\slshape    
We have seen that computer programming is an art, \\ 
because it applies accumulated knowledge to the world, \\ 
because it requires skill and ingenuity, and especially \\
because it produces objects of beauty.} \\ \medskip
--- \defcitealias{knuth:1974}{Donald E. Knuth}\citetalias{knuth:1974} \citep{knuth:1974}
\end{flushright}



\bigskip

%----------------------------------------------------------------------------------------

\begingroup

\let\clearpage\relax
\let\cleardoublepage\relax
\let\cleardoublepage\relax
%\addchaptertocentry{Remerciements} 
\addcontentsline{toc}{chapter}{Remerciements}
\chapter*{Remerciements}

Arrivant à l’aboutissement de ma tâche, je me trouve dans l’obligation respectueuse de devoir présenter mes chaleureux remerciements et témoignage de ma gratitude à tous ceux qui ont contribué aimablement et avec patience à l’élaboration de ce mémoire.

Tout d’abord, je tiens à remercier profondément les membres du jury qui m’ont fait l’honneur de juger mon travail. Merci à...

Je tiens à remercier toutes les personnes qui ont contribué au succès de ... et qui m'ont aidé lors de la rédaction de ce memoire. à ceux qui m’ont beaucoup appris au cours de la redaction, et même à ceux qui ont eu la gentillesse de faire de ce ... un moment très profitable.\\

Enfin, je tiens à remercier toutes les personnes qui m'ont conseillé et relu lors de la rédaction de ce rapport de stage : ma famille, mes ami(e)s et camarade de promotion.

\endgroup % Acknowledgements page

\pagestyle{scrheadings} % Show chapter titles as headings

\cleardoublepage% Table of Contents - List of Tables/Figures/Listings and Acronyms

\refstepcounter{dummy}

\pdfbookmark[1]{\contentsname}{tableofcontents} % Bookmark name visible in a PDF viewer

\setcounter{tocdepth}{2} % Depth of sections to include in the table of contents - currently up to subsections

\setcounter{secnumdepth}{3} % Depth of sections to number in the text itself - currently up to subsubsections

\manualmark
\markboth{\spacedlowsmallcaps{\contentsname}}{\spacedlowsmallcaps{\contentsname}}
\tableofcontents 
\automark[section]{chapter}
\renewcommand{\chaptermark}[1]{\markboth{\spacedlowsmallcaps{#1}}{\spacedlowsmallcaps{#1}}}
\renewcommand{\sectionmark}[1]{\markright{\thesection\enspace\spacedlowsmallcaps{#1}}}

\clearpage

\begingroup 
\let\clearpage\relax
\let\cleardoublepage\relax
\let\cleardoublepage\relax

%----------------------------------------------------------------------------------------
%	List of Figures
%----------------------------------------------------------------------------------------

\refstepcounter{dummy}
%\addcontentsline{toc}{chapter}{\listfigurename} % Uncomment if you would like the list of figures to appear in the table of contents
\pdfbookmark[1]{\listfigurename}{lof} % Bookmark name visible in a PDF viewer


%#
\listoffigures


\vspace{8ex}
\newpage

%----------------------------------------------------------------------------------------
%	List of Tables
%----------------------------------------------------------------------------------------

%\refstepcounter{dummy}
%\addcontentsline{toc}{chapter}{\listtablename} % Uncomment if you would like the list of tables to appear in the table of contents
%\pdfbookmark[1]{\listtablename}{lot} % Bookmark name visible in a PDF viewer

%\listoftables
        
%\vspace{8ex}
%\newpage
    
%----------------------------------------------------------------------------------------
%	List of Listings
%---------------------------------------------------------------------------------------- 

%\refstepcounter{dummy}
%\addcontentsline{toc}{chapter}{\lstlistlistingname} % Uncomment if you would like the list of listings to appear in the table of contents
%\pdfbookmark[1]{\lstlistlistingname}{lol} % Bookmark name visible in a PDF viewer

%\lstlistoflistings 

%\vspace{8ex}
%\newpage
       
%----------------------------------------------------------------------------------------
%	Acronyms
%----------------------------------------------------------------------------------------

\refstepcounter{dummy}
%\addcontentsline{toc}{chapter}{Acronyms} % Uncomment if you would like the acronyms to appear in the table of contents
\pdfbookmark[1]{Acronyms}{acronyms} % Bookmark name visible in a PDF viewer

\markboth{\spacedlowsmallcaps{Acronyms}}{\spacedlowsmallcaps{Acronyms}}

\chapter*{Liste des acronymes}

\begin{acronym}[UML]
	\acro{ML}{Machine Learning}
	\acro{CV}{Computer Vision}
	
	\acro{OCR}{Optical character recognition}
	\acro{ANPR}{Automatic number-plate recognition}
	\acro{ALPR}{Automatic license plate recognition}
	
	\acro{GD}{Gradient Descent}
	\acro{SGD}{Stochastic Gradient Descent}
	\acro{ADALINE}{ADAptative LInear NEuron }
	
	
	\acro{API}{Application Programming Interface}
	\acro{UML}{Unified Modeling Language}
\end{acronym} 
                   
\endgroup


%----------------------------------------------------------------------------------------
%	Symbol
%----------------------------------------------------------------------------------------
\refstepcounter{dummy}
%\addcontentsline{toc}{chapter}{Acronyms} % Uncomment if you would like the acronyms to appear in the table of contents
\pdfbookmark[1]{Notions}{notions} % Bookmark name visible in a PDF viewer

\markboth{\spacedlowsmallcaps{Notions}}{\spacedlowsmallcaps{Notions}}

\chapter*{Notions}

\begin{tabular}{ll}
	$\mathbb{N} $ & Ensemble des entiers naturels\\
	%$\mathbb{R} $ &  \\
	$\mathbb{R}^n $ & Ensemble des réels ou  Espace euclidien de dimension $n$ \\
	$\mathbb{B}^n = \{0,1\}^n $ & Espace booléen de dimension $n$\\
	$\mathcal{O}(\cdot) $ ou $ {\Omega}(\cdot) $ & L'ordre de grandeur maximal de complexité d'un algorithme \\
	%$\mathcal{O}(\cdot) $ & Le grand O de la notation de Landeau \\
	
	
	$x = \begin{pmatrix}
		x_1 \\ \vdots \\ x_n 
	\end{pmatrix} $ & Un vecteur\\ 
	$x =  (x_1, \vdots,  x_n)^T$ & Un vecteur \\ 
	$x^T =  (x_1, \vdots,  x_n)^T$ & Un vecteur transposé \\ 
	$ \langle xy\rangle = x^Ty$  & Le produit vectoriel \\ 
	$\parallel x \parallel $ & La norme du vecteur\\
	
	 
	$M^{-1}$ & La matrice inverse d'une matrice $M$\\
	$M^{T}$ & La matrice transposée \\
	
	$\frac{\partial}{\partial x}f(x,y) $ & La dérivée partielle par rapport à x de al fonction $f$ des deux variable $x$ et $y$ \\ 
	
	$\nabla_A J(A,B) $ & Le vecteur dérivé par rapport au vecteur $A$ de la fonctionnelle $J$ des deux vecteurs $A$ et $B$ \\ 
	
	$ $ & \\
	
	 $ $ & \textbf{  \ \ \ \ \ \ \ \ \textsc{Les éléments en apprentissage}} \\
	
	$ \mathcal{S} $ & L'échantillon d'apprentissage (un ensemble ou une suite d'exemple)  \\ 
	$ \hat{y}$ & valeurs prédite après l'entrainement d'un modèle d'apprentissage automatique \\ 
	$ \mathcal{H} $ & \\
	$ h \in \mathcal{H} $ & \\
	$y = h(x) $ & \\
	$ \ell(f(x),h(x)) $ & \\
	
	$ $ & \\
	
\end{tabular}
 % Contents, list of figures/tables/listings and acronyms

\cleardoublepage

\pagenumbering{arabic} % Arabic page numbering for thesis content (1, 2, 3, etc)
%\setcounter{page}{90} % Uncomment to manually start the page counter at an arbitrary value (for example if you wish to count the pre-content pages in the page count)

\cleardoublepage % Avoids problems with pdfbookmark
%------------------------------------------------
%#######################################################################################
%	CHAPTER - INTRO
%#######################################################################################

\cleardoublepage % Empty page before the start of the next part

%------------------------------------------------

\ctparttext{Ce stage à été effectué dans le but d'acquérir une connaissances sur l'application de l'informatique dans l'industrie. Nous allons parler de comment l'instrumentation et l'ingénierie système arrivent a faire la protection et sécurité des machines de production et le contrôle des procédés métallurgiques sur les différents sites de production.} % Text on the Part 2 page describing the content in Part 2

\textcolor{teal}{\part{Le cadre du stage (Instrumentation)}}

%#######################################################################################
%	CHAPTER 
%#######################################################################################



%----------------------------------------------------------------------------------------
%	THESIS CONTENT - APPENDICES
%----------------------------------------------------------------------------------------

\appendix

\ctparttext{Quelques programmes (code et schemas) pour les PLC et du DCS. Et Quelques images des diffentes sections}

\part{Annexes et Bibliographie} % New part of the thesis for the appendix



%% Appendix A

\chapter{Appendix Test}

%----------------------------------------------------------------------------------------

\lipsum[13-14]

%----------------------------------------------------------------------------------------

\section{Appendix Section Test}
\lipsum[15]

\graffito{More dummy text}
\lipsum[16]

%----------------------------------------------------------------------------------------

\section{Another Appendix Section Test}
\lipsum[17]

\begin{table}
\myfloatalign
\begin{tabularx}{\textwidth}{Xll} \toprule
\tableheadline{labitur bonorum pri no} & \tableheadline{que vista}
& \tableheadline{human} \\ \midrule
fastidii ea ius & germano &  demonstratea \\
suscipit instructior & titulo & personas \\
\midrule
quaestio philosophia & facto & demonstrated \\
\bottomrule
\end{tabularx}
\caption[Autem usu id]{Autem usu id.}
\label{tab:moreexample}
\end{table}

\lipsum[18]

There is also a useless Pascal listing below: \autoref{lst:useless}.

\begin{lstlisting}[float=b,language=Pascal,frame=tb,caption={A floating example (\texttt{listings} manual)},label=lst:useless]
for i:=maxint downto 0 do
begin
{ do nothing }
end;
\end{lstlisting} % Appendix A
%% Appendix X

\chapter{Appendix Title}

%----------------------------------------------------------------------------------------

% Content begins here % Appendix B - empty template

%----------------------------------------------------------------------------------------
%	POST-CONTENT THESIS PAGES
%----------------------------------------------------------------------------------------
%\bibliographystyle{plain}
%\addcontentsline{toc}{chapter}{\tocEntry{\bibname}}
%\bibliography{Bibliography}
\cleardoublepage% Bibliography

\label{app:bibliography} % Reference the bibliography elsewhere with \autoref{app:bibliography}

\manualmark % Work-around to have small caps also here in the headline
\markboth{\spacedlowsmallcaps{\bibname}}{\spacedlowsmallcaps{\bibname}} % Work-around to have small caps also
%\phantomsection
\refstepcounter{dummy}

\addtocontents{toc}{\protect\vspace{\beforebibskip}} % Place the bibliography slightly below the rest of the document content in the table of contents
\addcontentsline{toc}{chapter}{\tocEntry{\bibname}}

\printbibliography

 % Bibliography

%\cleardoublepage% Declaration

\refstepcounter{dummy}
\pdfbookmark[0]{Declaration}{declaration} % Bookmark name visible in a PDF viewer

\chapter*{Declaration} % Declaration section text

\thispagestyle{empty}
 

Ce rapport de stage effectué à Kamoto Copper Company a été rédigé par\\ \textbf{TSHELEKA KAJILA Hassan}, étudiant de l'\textsc{Université Nouveaux Horizons}, conformément aux exigences du diplôme de Licencié en Sciences informatique, département : \textbf{Calcul Scientifique}

\bigskip
 
\noindent\textit{\myLocation, \myTime}

\smallskip

\begin{flushright}
\begin{tabular}{m{5cm}}
\\ \hline
\centering\myName \\
\end{tabular}
\end{flushright}
 % Declaration

\cleardoublepage% Colophon (a brief description of publication or production notes relevant to the edition)

\pagestyle{empty}

\hfill

\vfill

\pdfbookmark[0]{Colophon}{colophon}

\section*{Colophon}
Cette étude a été très enrichissant pour moi, car il m'a permis de découvrir le domaine du Machine Learning, apprentissage supervisé, ses acteurs. Elle m'a permis de participer concrètement à ses enjeux au travers mes missions en apprentissage supervisé et la vision par ordinateur. Je préfère ainsi m'orienter vers un domaine lié à ma mission en Calcul Scientifique et le Data Science.



%This document was typeset using the typographical look-and-feel \texttt{classicthesis} developed by Andr\'e Miede. The style was inspired by Robert Bringhurst's seminal book on typography ``\emph{The Elements of Typographic Style}''. \texttt{classicthesis} is available for both \LaTeX\ and \mLyX: 

\begin{center}
\url{https://www.katangamining.com}
\end{center}
{Ce mémoire pour le travail de fin de cycle \textsf{$2021-2021$} qui traite la thématique du \texttt{Machine Learning \& Computer Vision} a été rédigé par {TSHELEKA KAJILA Hassan}, étudiant de l'\textsc{Université Nouveaux Horizons}, conformément aux exigences du diplôme de Licencié en Sciences informatique, département : \emph{Calcul Scientifique}.}


\begin{center}
\url{https://www.unhorizons.org}
\end{center}
 
\bigskip

\noindent\finalVersionString % Colophon

%----------------------------------------------------------------------------------------

\end{document}
